% Point d'avancement tuteur, 21 aout 2026. DIX pages : titre, rappel, et huit slides de fond.
%
% CE FICHIER REMPLACE UNE VERSION ECRITE LE 5 AOUT, donc AVANT les points du 7 et du 14. Cette
% version-la presentait le biais de narration, les quatre bilans reels et l'invariance de g,
% c'est-a-dire du travail ANTERIEUR au deck du 14, tout en annoncant en premiere slide qu'elle
% le suivait. Elle portait en outre les valeurs MACSF corrigees depuis (3 234 d'actifs, 292 de
% SCR, 41,3 % de part) et l'argument du levier a 13 % qui vaut 9 %. Le remplacement est donc a
% la fois une mise en ordre chronologique et la correction de valeurs qu'on sait fausses, seule
% exception au gel d'un deck.
%
% CONTINUITE DU 14/08, VERIFIEE POINT PAR POINT :
%   - le 14/08 publiait la table complete des quatre canaux -> la section A tranche LEQUEL des
%     quatre protocoles porte le chiffre de tete, question que le deck precedent laissait
%     ouverte sans le dire ;
%   - le 14/08 empruntait cinq notions au preprint de cascade climatique -> la section B TESTE
%     celle qui promettait le plus (le coin superieur) et la section C ecrit le desaccord ;
%   - le 14/08 laissait la posture reportee en attente -> la section D la tranche, avec trois
%     autres decisions.
% FIL : cette semaine n'a presque rien mesure de neuf, elle a TRANCHE.
%
% DEUX SUJETS SONT HORS DU DECK, SUR DECISION DE KELIAN. NE PAS LES REMETTRE.
%   1. LA CHARTE GRAPHIQUE. Le passage des 92 scripts a une source unique de couleurs et les 49
%      figures sur 50 ne sont pas un resultat de modelisation : c'est de l'outillage, et cela
%      n'a pas sa place dans une heure avec le tuteur.
%   2. TOUT CE QUI RELEVE DE L'ECRITURE DU MEMOIRE. Donc pas de comptes de pages, pas de taux du
%      harnais, pas de defauts de citation, pas de renvois casses, pas de couverture de deck.
%      Ce travail existe et il est reel, mais il decrit la FABRICATION du document et non ses
%      resultats. Il est trace dans CLAUDE.md et dans PLAN_MEMOIRE.md.
%   Consequence sur les sources : les lignes « Sources : scripts... » restent, elles sont la
%   pour le harnais et elles ne coutent rien a l'oral, mais AUCUNE slide ne parle du harnais.
%
% REGLE DE DENSITE, ET C'EST UNE CONSIGNE : le deck est un SUPPORT DE PAROLE, pas un document.
% Une slide porte le titre qui enonce le resultat, trois a cinq lignes courtes, et les chiffres.
% Le raisonnement passe a l'oral. Une redaction anterieure faisait des paragraphes de quinze
% lignes par bloc : elle debordait sur huit slides et elle se lisait au lieu de s'ecouter.
% NE PAS LA REDENSIFIER. Si une idee ne tient pas en une ligne, elle est deja dans le memoire.
%
% VOCABULAIRE : point tuteur, donc tutoiement et raccourcis du projet employes sans les definir
% (lecture a quatre canaux, marginal de fermeture, Shapley, graine, pave, IC90, plug-in). Hugo
% les connait, les expliquer sur une slide coute une slide.
%
% AUCUNE FIGURE, comme le deck du 14 aout. Motif : la semaine a produit des decisions et des
% tables, pas de nouvelles courbes, et les trois figures de la version remplacee (Z22, S22, S23)
% portent du travail qui n'est pas celui de cette semaine. Si une figure est souhaitee, S10
% (SCR par etat) et S24 (interaction des canaux) n'ont jamais servi dans un deck.
%
% PUCES : le template beamer est force en cercle. Le babel francais met un tiret cadratin en
% puce d'itemize par defaut, et les tirets cadratins sont proscrits dans les livrables. Le
% controle se fait SUR LE PDF, jamais sur la source.
%
% HARNAIS : chaque section porte sa ligne de sources ; la page de titre et la slide de rappel
% sont declarees hors section, motif ecrit. Une declaration n'est honoree que si la section ne
% cite AUCUN script, et elle doit etre posee APRES le \section* : posee avant, elle tombe dans
% le bloc precedent. Vu ici meme.
%
% CONTROLES PASSES : 0 Overfull vbox, 0 Overfull hbox, 0 tiret cadratin releve sur le PDF,
% 0 « ?? ». Harnais : 100 % de confirmation, HORS CONTROLE NON DECLARE = 0.
\documentclass[11pt]{beamer}
\usetheme{Madrid}
\usecolortheme{seahorse}
\usepackage[utf8]{inputenc}
\usepackage[T1]{fontenc}
\usepackage[french]{babel}
\usepackage{booktabs}
\usepackage{amsmath,amssymb}
\usepackage{eurosym}
\graphicspath{{../vasicek_lab/figures/}{../cascade_qualitative/figures/}{../vasicek_lab/notes/}}
\setbeamertemplate{navigation symbols}{}
\setbeamertemplate{caption}{\raggedright\insertcaption\par}
\setbeamertemplate{itemize items}[circle]
\setbeamerfont{block title}{size=\small}

\title[Point d'avancement]{Mémoire SCR DORA : point d'avancement}
\subtitle{La semaine où il a fallu trancher}
\author[K. Kaddouri]{Kélian Kaddouri}
\institute[ENSAE / Nexialog]{ENSAE / Nexialog Consulting}
\date[21/08/2026]{21 août 2026}

\begin{document}

% HARNAIS-HORS-SECTION: page de titre. Millesime et date, aucun resultat.
\begin{frame}
\titlepage
\end{frame}

% =====================================================================
\section*{Rappel : le point précédent}
% HARNAIS-HORS-SECTION: slide de rappel. Le seul nombre est la DATE du point precedent. Les
% grandeurs rappelees le sont par leur NOM, precisement pour que ce rappel n'ait pas a etre
% reverifie : elles ont ete confirmees dans le deck du 14.

\begin{frame}{Depuis le point du 14 août}
\large
Le 14 : les quatre canaux, Shapley et Euler, portage et plancher, les six postures.
\vspace{5mm}

\normalsize
\textbf{Cette semaine : quatre résultats et une limite qui reste.}
\vspace{4mm}
\begin{itemize}\itemsep5pt
  \item[\textbf{A}] Quatre protocoles chiffraient l'écart. Un seul est publié
  \item[\textbf{B}] Le coin supérieur : testé, il ne transporte qu'à moitié
  \item[\textbf{C}] La brique la plus lourde est la queue, pas la propagation
  \item[\textbf{D}] Quatre décisions de modélisation
  \item[\textbf{E}] La direction de $W$, seul point non mesuré
\end{itemize}
\end{frame}

% =====================================================================
\section*{A. Le chiffre de tête}

\begin{frame}{Quatre protocoles chiffraient le même écart}
\begin{center}\small
\begin{tabular}{@{}l c r r r r@{}}
\toprule
\textbf{Lecture} & \textbf{can.} & \textbf{conf.} & \textbf{non conf.} & \textbf{fact.} &
\textbf{écart} \\
\midrule
Score seul, $g$ figé NC & $1$ & $7\,987$ & $16\,847$ & $2{,}11$ & $8\,860$ \\
Fréquence $+$ propagation (B) & $2$ & $6\,664$ & $15\,074$ & $2{,}26$ & $8\,410$ \\
$+$ détection (C) & $3$ & $5\,932$ & $16\,595$ & $2{,}80$ & $10\,663$ \\
\textbf{$+$ accumulation P4} & $\mathbf{4}$ & $\mathbf{6\,049}$ & $\mathbf{20\,188}$ &
$\mathbf{3{,}34}$ & $\mathbf{14\,139}$ \\
\bottomrule
\end{tabular}
\end{center}
\vspace{5mm}
\begin{itemize}\itemsep5pt
  \item Aucun n'est faux : ils ne relâchent pas le même nombre de canaux
  \item \alert{Il fallait donc dire lequel je publie, et sous quelle définition}
\end{itemize}
\vspace{3mm}
{\scriptsize Sources : scripts~\texttt{16}, \texttt{25}, \texttt{43} et~\texttt{67}.}
\end{frame}

% =====================================================================
\begin{frame}{Je publie la lecture à quatre canaux}
\vspace{3mm}
\begin{center}
{\Large\alert{$\mathbf{6\,049 \to 20\,188}$~\textbf{M\euro}}}\\[4pt]
{\large facteur $3{,}34$ \quad écart $14\,139$~M\euro}
\end{center}
\vspace{4mm}
\begin{itemize}\itemsep5pt
  \item Seule lecture où l'état conforme est conforme sur \textbf{tous} les canaux
  \item Ailleurs, c'est une remédiation partielle qu'on appelle conformité
  \item Ce qui débloque : détection validée à $2\,928 \pm 343$ en Shapley, signe résolu sur
        seize graines
  \item Les trois autres restent publiées, comme remédiations partielles
  \item \alert{Comparer le facteur, pas l'écart en euros : la base change}
\end{itemize}
\vspace{3mm}
{\scriptsize Sources : scripts~\texttt{43}, \texttt{67} et~\texttt{68}.}
\end{frame}

% =====================================================================
\section*{B. Le théorème emprunté, testé}

\begin{frame}{Le coin supérieur : ce qui tient}
\large
Si l'état non conforme majore la perte sur tout le pavé, un test de résistance sur les quatre
canaux devient \textbf{une seule évaluation}.
\vspace{6mm}

\normalsize
\begin{itemize}\itemsep6pt
  \item Seize configurations, $65$ paires emboîtées
  \item \textbf{Zéro violation}, graine par graine et pas sur la moyenne
  \item Les deux coins redonnent $6\,049$ et $20\,188$, sinon je démontrais autre chose
\end{itemize}
\vspace{5mm}
\alert{\textbf{Au sens du quantile, c'est démontré.}}
\vspace{3mm}

{\scriptsize Sources : script~\texttt{88}.}
\end{frame}

% =====================================================================
\begin{frame}{Le coin supérieur : ce qui ne tient pas}
\small
Leur théorème est \textbf{trajectoriel}. À aléas communs, la perte d'une année baisse :
\vspace{-1mm}
\begin{center}\small
\begin{tabular}{@{}l r r@{}}
\toprule
\textbf{Canal relâché seul} & \textbf{années en baisse} & \textbf{part} \\
\midrule
Propagation & $6\,470$ & $16{,}18\,\%$ \\
\textbf{Détection} & $\mathbf{12\,275}$ & $\mathbf{30{,}69\,\%}$ \\
Accumulation tiers & $408$ & $1{,}02\,\%$ \\
\bottomrule
\end{tabular}
\end{center}
\vspace{4mm}
\begin{itemize}\itemsep6pt
  \item La détection viole le plus, et \textbf{pas à cause de la cascade} : $p_u$ agit sur la
        transformation de sévérité, pas sur une table
  \item \alert{Un coin peut être infaisable} : $2{,}0$ contre $1{,}2$ sur leur contre-exemple
  \item Deux canaux sur quatre sont des bornes posées, donc c'est un majorant sur un pavé
        déclaré, pas une entité
\end{itemize}
\vspace{3mm}
{\scriptsize Sources : script~\texttt{88}.}
\end{frame}

% =====================================================================
\section*{C. La brique la plus lourde}

\begin{frame}{Ce qui porte le capital, c'est la queue, pas la cascade}
\small
J'enlève les briques une par une et je regarde ce qu'il reste du capital.
\vspace{3mm}

\begin{itemize}\itemsep5pt
  \item Retirer la \textbf{queue lourde} : $-76{,}2\,\%$
  \item Retirer la \textbf{propagation} : $-20{,}3\,\%$
  \item \alert{Facteur $3{,}7$, même conclusion que le tornado par un chemin indépendant}
  \item La forme de queue est \textbf{héritée de la sévérité} : $\mathrm{CTE}/\mathrm{VaR}$
        passe de $2{,}075$ à $2{,}061$ seulement
\end{itemize}
\vspace{4mm}
\begin{block}{Le préprint conclut l'inverse, et la cause est dans son texte}
Leur brique la plus lourde est la propagation, $-53{,}9\,\%$. Mais leur sévérité est
\textbf{bornée}, donc sans indice de queue, et leur ablation ne contient aucune brique de queue.
Les deux classements sont corrects chacun dans son modèle.
\end{block}
\vspace{2mm}
{\scriptsize Sources : scripts~\texttt{63} et~\texttt{81}.}
\end{frame}

% =====================================================================
\section*{D. Quatre décisions}

\begin{frame}{Posture reportée : le plug-in avec sa bande}
\begin{itemize}\itemsep6pt
  \item Motif réglementaire d'abord : le \textsc{scr} est une VaR à $99{,}5\,\%$, pas un
        suprémum sur une famille de lois
  \item La prédictive ne déplace pas le point au niveau réglementaire
  \item Les postures robustes sont les moins reproductibles des six
\end{itemize}
\vspace{5mm}
\begin{block}{Et l'identité que j'ai trouvée en répondant}
La \textbf{robuste à $95\,\%$ est la borne haute de l'IC$90$} que je publie déjà. Même loi
bootstrap, $1\,037$~M\euro{} des deux côtés.
\end{block}
\vspace{3mm}
\begin{itemize}\itemsep5pt
  \item Donc publier le plug-in \emph{avec sa bande} publie déjà le chiffre robuste
  \item \alert{Le débat porte sur le nom du capital, pas sur l'information disponible}
\end{itemize}
\vspace{3mm}
{\scriptsize Sources : scripts~\texttt{46} et~\texttt{47}.}
\end{frame}

% =====================================================================
\begin{frame}{Les trois autres}
\begin{center}\small
\begin{tabular}{@{}l l l@{}}
\toprule
\textbf{Point} & \textbf{Décision} & \textbf{Motif} \\
\midrule
Niveau de l'intervalle & $90\,\%$ & le manque est le même à $95$ \\
\textsc{cte} en diagnostic & déjà fait & avec le niveau équivalent \\
Ancrage de $g$ (\textsc{acpr}, \textsc{eiopa}) & non & pas la bonne grandeur publiée \\
\bottomrule
\end{tabular}
\end{center}
\vspace{5mm}
\begin{itemize}\itemsep6pt
  \item Couverture réelle $86{,}8\,\%$ pour un nominal de $90$, $91{,}8\,\%$ pour $95$ :
        \alert{le manque est $-3{,}2$ points dans les deux cas}, donc relever le nominal ne
        corrige pas l'estimateur
  \item $\mathrm{CTE}_\beta$ égale la VaR réglementaire à $\beta = 97{,}73\,\%$ ; à $95\,\%$ elle
        vaut $5\,734$ contre $8\,374$, donc \textbf{moins} prudente que l'exigence
  \item Ces sources publient des attentes prudentielles, pas des probabilités de propagation :
        un ancrage aurait l'apparence d'un calibrage sans en être un
\end{itemize}
\vspace{3mm}
{\scriptsize Sources : scripts~\texttt{72} et~\texttt{73}.}
\end{frame}

% =====================================================================
\section*{E. La direction de W}

\begin{frame}{Le seul point non mesuré : la direction de \texorpdfstring{$W$}{W}}
\small
La co-occurrence est identifiée par la donnée. La \textbf{direction}, non : je la code moi-même
sur les post-mortems, donc elle repose sur \textbf{un seul codeur}. C'est l'objection qu'un jury
formulera en premier.
\vspace{3mm}

\begin{itemize}\itemsep5pt
  \item Le kit est prêt : dix récits anonymisés, sans aucune de mes flèches
  \item Une heure de lecture, cent cases. Toi, ou un praticien hors du projet
  \item Le script ne fabrique rien sans donnée
\end{itemize}
\vspace{4mm}
\begin{block}{Ce que je peux déjà dire si personne ne code}
Je borne l'objection par l'adverse. En rendant $P_1$ muette, le capital d'entité bouge de
$41{,}9$~M\euro, soit \textbf{$0{,}81$ fois} la bande d'identification que je publie déjà.
\alert{Accorder au critique toute son objection coûte moins que l'ignorance directionnelle que
j'assume.} Ce qui bouge, en revanche, c'est le pilier à remédier en premier.
\end{block}
\vspace{2mm}
{\scriptsize Sources : scripts~\texttt{59} et~\texttt{64}.}
\end{frame}

\end{document}
